\begin{mistake}{2026-04-09}{48}

\begin{problemcontent}
设连续函数 \( P(x), Q(x) \) 以 \( T \) 为周期，\( y = y(x) \) 是方程 \( y' + P(x)y = Q(x) \) 在 \( (-\infty, +\infty) \) 上的解。证明：\( y(0) = y(T) \) 是 \( y(x) \) 以 \( T \) 为周期的充分必要条件。
\end{problemcontent}

\begin{solutioncontent}
\textbf{必要性}：若 \( y(x) \) 是以 \( T \) 为周期的周期函数，根据定义有对于任意实数 \( x \)，\( y(x+T) = y(x) \) 恒成立。令 \( x = 0 \)，即得 \( y(T) = y(0) \)。

\textbf{充分性}：已知 \( y(0) = y(T) \)，需证 \( y(x+T) \equiv y(x) \)。
构造函数 \( u(x) = y(x+T) \)。
因为 \( y(x) \) 满足原微分方程，所以将其自变量替换为 \( x+T \) 得到：
\[ y'(x+T) + P(x+T)y(x+T) = Q(x+T) \]
由于 \( P(x) \) 和 \( Q(x) \) 的周期均为 \( T \)，即 \( P(x+T) = P(x) \), \( Q(x+T) = Q(x) \)。
代入上式得：
\[ u'(x) + P(x)u(x) = Q(x) \]
这说明 \( u(x) \) 和 \( y(x) \) 满足完全相同的一阶线性微分方程。
再看初始条件：
\( u(0) = y(T) \)，由已知 \( y(T) = y(0) \)，故 \( u(0) = y(0) \)。
根据一阶常微分方程初值问题解的唯一性定理，满足相同方程和相同初值的解必定恒等，故在 \(( -\infty, +\infty )\) 上：
\[ u(x) \equiv y(x) \implies y(x+T) \equiv y(x) \]
即 \( y(x) \) 以 \( T \) 为周期。充分性得证。
\end{solutioncontent}

\mistakeknowledge{
\begin{itemize}
 \item \textbf{核心定理}：常微分方程初值问题解的存在唯一性定理；周期函数的平移代数定义。
 \item \textbf{易错点}：证明充分性时思路受阻，企图利用一阶线性非齐次方程的通解积分公式直接硬算验证，过程繁杂且极易因为定积分上下限处理不当而出错。
 \item \textbf{关联知识}："平移构造法"是处理ODE周期性及对称性相关证明题最为通用和优雅的手法。
\end{itemize}}

\end{mistake}
